%% kapa: exam=1 round=35 year=2024 period=2 form=A subject=accounting subject_ko=회계학 seq=12 number=52
%% source: Q-Net 감정평가사 기출문제 (cst00309) artlSeq=5213229 file=raw/1008/5213229/2교시 문제지 원본.pdf
%% conversion: kordoc 4.24.51 -> tools/gen1.py (auto); review=proofread-sonnet; flags=none
\begin{question}{52}
(주)감평은 20×1년 1월 1일 다음과 같은 조건의 비분리형 신주인수권부사채를 액면발행하였다.
\begin{framedtext}
○ 액면금액: ￦1,000 \par
○ 표시이자율: 연 5\% \par
○ 사채발행시 신주인수권이 부여되지 않은 일반사채의 시장이자율: 연 12\% \par
○ 이자지급일: 매년 12월 31일 \par
○ 행사가격: 1주당 ￦200 \par
○ 발행주식의 액면금액: 1주당 ￦100 \par
○ 만기상환일: 20×3년 12월 31일 \par
○ 상환조건: 신주인수권 미행사시 상환기일에 액면금액의 113.5\%를 일시상환
\end{framedtext}
20×2년 초 상기 신주인수권의 60\%가 행사되어 3주가 발행되었다. 20×2년 초 상기 신주인수권의 행사로 인해 증가하는 (주)감평의 주식발행초과금은? (단, 신주인수권 행사시 신주인수권대가는 주식발행초과금으로 대체한다. 화폐금액은 소수점 첫째자리에서 반올림하며, 단수차이로 인한 오차는 가장 근사치를 선택한다.)
\par{\small\centering\begin{tabular}{|c|c|c|c|c|}\hline
\multirow{2}{*}{기간} & \multicolumn{2}{c|}{단일금액 ￦1의 현재가치} & \multicolumn{2}{c|}{정상연금 ￦1의 현재가치} \tabularnewline\cline{2-5}
 & 5\% & 12\% & 5\% & 12\% \tabularnewline\hline
1 & 0.9524 & 0.8928 & 0.9524 & 0.8928 \tabularnewline\hline
2 & 0.9070 & 0.7972 & 1.8594 & 1.6900 \tabularnewline\hline
3 & 0.8638 & 0.7118 & 2.7232 & 2.4018 \tabularnewline\hline
\end{tabular}\par}
\begin{choices}
  \item ￦308
  \item ￦335
  \item ￦365
  \item ￦408
  \item ￦435
\end{choices}
\end{question}
\answer{4}
